% =====================================================================
%  Stage 2, раздел 1.
%  Сравнительный анализ архитектур нейронных сетей для NIDS.
%  Опубликованные результаты применения архитектур
%  с явным указанием первоисточников.
% =====================================================================

\section{Сравнительный анализ архитектур нейронных сетей для обнаружения вторжений}

\subsection{Принципы и методология сравнения}

Цель раздела --- сопоставить архитектуры нейронных сетей, рассмотренные в
теоретической части (Этап~1, раздел~1.4), на основе экспериментальных
результатов, опубликованных в рецензируемых источниках в применениях,
структурно эквивалентных задаче настоящей работы --- бинарной/мультиклассовой
детекции аномалий в flow-представлении сетевого трафика.

Для каждой архитектуры приводятся:
\begin{itemize}
    \item первоисточник (научная работа, в которой опубликован результат);
    \item использованный датасет;
    \item постановка задачи (\emph{binary} либо \emph{multiclass});
    \item конкретное количественное значение метрики, опубликованное в работе.
\end{itemize}

\textbf{Методологическое замечание.} Прямое количественное сравнение
результатов из разных источников некорректно по следующим причинам.
\begin{enumerate}
    \item Используются разные датасеты (NSL-KDD, KDD Cup 99, UNSW-NB15,
    CICIDS2017) с принципиально различающимися моделями угроз и распределениями
    признаков~\cite{moustafa2016eval, sommer2010outside, sharafaldin2018toward}.
    \item Используются различные схемы препроцессинга, селекции признаков и
    разделения выборок~\cite{khraisat2019survey, ferrag2020dlcyber}.
    \item Метрики (\emph{Accuracy, Detection Rate, F1, Precision, Recall, AUC})
    интерпретируются в разных задачах различно; в условиях сильного
    дисбаланса классов \emph{Accuracy} даёт оптимистическую
    оценку~\cite{khraisat2019survey, liu2019survey}.
\end{enumerate}

Поэтому опубликованные результаты используются здесь как качественные
индикаторы пригодности архитектуры в NIDS-контексте и для теоретического
обоснования предположения о работоспособности модели на UNSW-NB15. Прямое
количественное сравнение моделей-кандидатов в равных условиях будет
выполнено в собственном эксперименте Этапа~4 настоящей работы.

\subsection{Расширенная теоретическая база сравниваемых архитектур}

В раздел~1 Этапа~1 вошло базовое описание архитектур; здесь приводятся
дополнительные характеристики, существенные для NIDS-применения.

\subsubsection{Полносвязные сети (DNN/MLP)}

В контексте flow-данных DNN рассматривается как нелинейный классификатор
над агрегированными признаками отдельного flow-вектора. Стандартная
композиция включает 2--5 скрытых слоёв с ReLU-активациями и dropout-регуляризацией.
Глубина свыше пяти слоёв в задачах NIDS, по данным
обзоров~\cite{vinayakumar2019dl, ferrag2020dlcyber, ahmad2021nids}, не
приводит к устойчивому улучшению качества.

\subsubsection{1D-CNN над окнами flow-векторов}

Приложение свёрточных слоёв в NIDS интерпретируется как извлечение
\emph{мотивов} (\eng{motifs}) во времени или в признаковой строке.
Стандартная конфигурация: один--три блока \texttt{Conv1D-BatchNorm-ReLU}
с пуллингом, размер ядра \(K \in \{3, 5, 7\}\). При обработке
последовательности окон используется ширина окна
\(W \in \{8, 16, 32, 64\}\), что согласуется с эксплуатационными
характеристиками IPFIX-коллекторов~\cite{rfc7011}.

\subsubsection{Рекуррентные сети: LSTM, GRU, BiLSTM}

В flow-представлении один временной шаг соответствует одному соединению или
агрегированному окну. LSTM/GRU моделируют долгосрочные зависимости между
соединениями, что востребовано для атак с нерегулярной временной структурой
(slow brute-force, low-and-slow exfiltration). BiLSTM использует контекст в
обоих направлениях, что даёт прирост качества за счёт повышенной задержки.

\subsubsection{Гибрид CNN--LSTM}

Иерархия \texttt{Conv1D $\to$ Pooling $\to$ LSTM $\to$ FC} извлекает локальные
паттерны (CNN) и моделирует их временную динамику (LSTM). Гибрид является
устоявшейся и активно публикуемой
архитектурой~\cite{vinayakumar2019dl, ferrag2020dlcyber, ahmad2021nids}.

\subsubsection{TCN}

\emph{Temporal Convolutional Network}~\cite{bai2018tcn} реализует
последовательное моделирование на основе причинных дилатированных свёрток.
В NIDS-публикациях TCN встречается реже LSTM/CNN-LSTM; преимущество ---
параллелизуемость; ограничение --- фиксированное рецептивное поле.

\subsubsection{Transformer-encoder}

Применение Transformer к flow-последовательностям требует позиционного
кодирования и достаточного объёма данных для устойчивого обучения. В
NIDS-публикациях Transformer применяется реже CNN-LSTM;
обзоры~\cite{ferrag2020dlcyber, ahmad2021nids} отмечают сопоставимое или
несколько лучшее качество, но повышенную дисперсию по seed.

\subsubsection{Autoencoder и VAE}

Автоэнкодер обучается на нормальном трафике и идентифицирует аномалии
через ошибку реконструкции; такой режим (\emph{semi-supervised}) ценен
тем, что не требует разметки атак. Эталонная online-реализация для
NIDS --- Kitsune~\cite{mirsky2018kitsune} --- использует ансамбль
автоэнкодеров с инкрементальным обучением.

\subsection{Опубликованные результаты применения DNN/MLP}

В \tabref{tab:dnn-results} сведены опубликованные результаты применения
DNN/MLP-моделей к NIDS-датасетам. Источниками являются рецензируемые работы
по обнаружению вторжений на основе глубокого обучения.

\begin{table}[H]
\centering
\caption{Опубликованные результаты применения DNN/MLP в NIDS}
\label{tab:dnn-results}
\renewcommand{\arraystretch}{1.25}
\setlength{\tabcolsep}{4pt}
\small
\begin{tabularx}{\linewidth}{@{}lXlXX@{}}
\toprule
\textbf{Источник} & \textbf{Архитектура} & \textbf{Датасет} &
\textbf{Постановка} & \textbf{Опубликованный результат} \\
\midrule
Vinayakumar и соавт.~\cite{vinayakumar2019dl} &
DNN, 1--5 скрытых слоёв (\eng{DNN1\dots DNN5}) &
NSL-KDD, KDD Cup 99, UNSW-NB15, CICIDS2017, WSN-DS &
binary и multiclass &
В зависимости от датасета и числа слоёв; авторы рекомендуют DNN из 3--5 слоёв
как универсальную конфигурацию по совокупности accuracy, precision, recall \\
\addlinespace
Sharafaldin и соавт.~\cite{sharafaldin2018toward} &
MLP &
CICIDS2017 &
multiclass &
Precision $\approx 0{,}77$; Recall $\approx 0{,}83$
(уступает RF/ID3, опубликованным в той же работе) \\
\addlinespace
Moustafa и Slay~\cite{moustafa2016eval} &
ANN &
UNSW-NB15 &
binary &
Accuracy $\approx 81{,}3\%$; FAR $\approx 18{,}4\%$ \\
\addlinespace
Khraisat и соавт.~\cite{khraisat2019survey} (обзор) &
DNN &
NSL-KDD, UNSW-NB15 &
binary, multiclass &
Сводные диапазоны Accuracy
$80\text{--}95\%$ для современных DL-моделей; результаты сильно
зависят от препроцессинга \\
\bottomrule
\end{tabularx}
\end{table}

Из \tabref{tab:dnn-results} видно, что DNN/MLP --- адекватная нижняя
граница DL-моделей в NIDS, но устойчиво уступает архитектурам, способным
учитывать локальные и временные структуры (см. далее).

\subsection{Опубликованные результаты применения 1D-CNN}

В \tabref{tab:cnn-results} приведены опубликованные результаты применения
свёрточных архитектур.

\begin{table}[H]
\centering
\caption{Опубликованные результаты применения 1D-CNN в NIDS}
\label{tab:cnn-results}
\renewcommand{\arraystretch}{1.25}
\setlength{\tabcolsep}{4pt}
\small
\begin{tabularx}{\linewidth}{@{}lXlXX@{}}
\toprule
\textbf{Источник} & \textbf{Архитектура} & \textbf{Датасет} &
\textbf{Постановка} & \textbf{Опубликованный результат} \\
\midrule
Vinayakumar и соавт.~\cite{vinayakumar2019dl} &
1D-CNN с пуллингом &
NSL-KDD, UNSW-NB15, KDD Cup 99, CICIDS2017 &
binary и multiclass &
Стабильное превосходство над DNN/MLP по metric F1 на тех же датасетах,
по данным авторов; усиление эффекта при увеличении длины окна \\
\addlinespace
Ferrag и соавт.~\cite{ferrag2020dlcyber} (обзор и собственные эксперименты) &
1D-CNN &
CSE-CIC-IDS2018 и Bot-IoT &
multiclass &
1D-CNN включён в семёрку наиболее перспективных архитектур по сводной
оценке F1 на крупных современных датасетах \\
\addlinespace
Liu и Lang~\cite{liu2019survey} (обзор) &
CNN различных конфигураций &
KDD-семейство, UNSW-NB15 &
binary, multiclass &
Опубликованные значения accuracy в применениях CNN $\geq 90\%$ для большинства
рассмотренных работ; качество выше при использовании CNN-блоков на окнах
последовательных flow \\
\bottomrule
\end{tabularx}
\end{table}

\subsection{Опубликованные результаты применения LSTM/GRU/BiLSTM}

В \tabref{tab:rnn-results} приведены опубликованные результаты для
рекуррентных архитектур.

\begin{table}[H]
\centering
\caption{Опубликованные результаты применения LSTM/GRU/BiLSTM в NIDS}
\label{tab:rnn-results}
\renewcommand{\arraystretch}{1.25}
\setlength{\tabcolsep}{4pt}
\small
\begin{tabularx}{\linewidth}{@{}lXlXX@{}}
\toprule
\textbf{Источник} & \textbf{Архитектура} & \textbf{Датасет} &
\textbf{Постановка} & \textbf{Опубликованный результат} \\
\midrule
Kim и соавт.~\cite{kim2016lstm} &
LSTM-RNN &
KDD Cup 99 &
binary &
Accuracy $\approx 96{,}93\%$; Detection Rate $\approx 98{,}88\%$;
False Alarm Rate $\approx 10{,}04\%$ \\
\addlinespace
Yin и соавт.~\cite{yin2017dlrnn} &
RNN (стандартная) &
NSL-KDD, KDDTest+ &
binary &
Accuracy $\approx 83{,}28\%$ (при 80 hidden units); существенное падение
на KDDTest-21 \\
\addlinespace
Yin и соавт.~\cite{yin2017dlrnn} &
RNN (стандартная) &
NSL-KDD, KDDTest+ &
multiclass &
Accuracy $\approx 81{,}29\%$ \\
\addlinespace
Vinayakumar и соавт.~\cite{vinayakumar2019dl} &
LSTM, GRU &
NSL-KDD, KDD Cup 99, UNSW-NB15 &
binary, multiclass &
LSTM и GRU сопоставимы по средней Accuracy и F1 на NSL-KDD/KDD Cup~99;
GRU стабильно дешевле по числу параметров; на UNSW-NB15 LSTM уступает
гибриду CNN-LSTM в той же работе \\
\addlinespace
Ahmad и соавт.~\cite{ahmad2021nids} (обзор) &
LSTM, BiLSTM &
NSL-KDD, UNSW-NB15, CICIDS2017 &
binary, multiclass &
BiLSTM устойчиво даёт прирост Accuracy/F1 относительно LSTM на NSL-KDD и
UNSW-NB15 за счёт двунаправленного контекста \\
\bottomrule
\end{tabularx}
\end{table}

Особо отметим методологические оговорки:
\begin{itemize}
    \item Высокая Accuracy и DR в~\cite{kim2016lstm} получены на
    KDD Cup 99, корпусе с известными статистическими
    дефектами~\cite{tavallaee2009nslkdd, sommer2010outside}; результат
    рассматривается как историческая точка отсчёта, а не как практически
    переносимая оценка.
    \item Yin и соавт.~\cite{yin2017dlrnn} явно отмечают разрыв качества
    между KDDTest+ и более сложным KDDTest-21, что подчёркивает
    зависимость опубликованных метрик от выборки тестирования.
\end{itemize}

\subsection{Опубликованные результаты применения CNN--LSTM}

В \tabref{tab:cnnlstm-results} сведены опубликованные результаты применения
гибрида CNN--LSTM в задачах NIDS.

\begin{table}[H]
\centering
\caption{Опубликованные результаты применения гибрида CNN--LSTM в NIDS}
\label{tab:cnnlstm-results}
\renewcommand{\arraystretch}{1.25}
\setlength{\tabcolsep}{4pt}
\small
\begin{tabularx}{\linewidth}{@{}lXlXX@{}}
\toprule
\textbf{Источник} & \textbf{Архитектура} & \textbf{Датасет} &
\textbf{Постановка} & \textbf{Опубликованный результат} \\
\midrule
Vinayakumar и соавт.~\cite{vinayakumar2019dl} &
CNN-LSTM &
NSL-KDD, KDD Cup 99, UNSW-NB15, CICIDS2017 &
binary, multiclass &
В семействе DL-моделей данной работы CNN-LSTM показал одни из
наибольших значений F1 и Accuracy на UNSW-NB15 и CICIDS2017
относительно DNN/CNN/LSTM той же работы \\
\addlinespace
Ferrag и соавт.~\cite{ferrag2020dlcyber} &
CNN-LSTM &
CSE-CIC-IDS2018, Bot-IoT, других &
multiclass &
CNN-LSTM фигурирует в перечне эталонных DL-моделей с устойчивыми F1 на
тестовых разделениях; авторы фиксируют преимущество в обнаружении атак с
выраженной временной структурой \\
\addlinespace
Liu и Lang~\cite{liu2019survey} (обзор) &
Гибридные CNN-RNN &
KDD-семейство, UNSW-NB15 &
binary, multiclass &
Гибридные подходы CNN-RNN наиболее устойчиво превосходят как чисто
свёрточные, так и чисто рекуррентные базы по metric F1 в большинстве
рассмотренных работ \\
\addlinespace
Ahmad и соавт.~\cite{ahmad2021nids} (обзор) &
CNN-LSTM, BiLSTM &
NSL-KDD, UNSW-NB15, CICIDS2017 &
binary, multiclass &
В обзоре отмечается, что CNN-LSTM и BiLSTM достигают наилучших Accuracy и F1
в большинстве экспериментов на рассматриваемых датасетах \\
\bottomrule
\end{tabularx}
\end{table}

\subsection{Опубликованные результаты применения автоэнкодеров и VAE}

В \tabref{tab:ae-results} сведены опубликованные результаты применения
автоэнкодерных архитектур в NIDS.

\begin{table}[H]
\centering
\caption{Опубликованные результаты применения AE/VAE в NIDS}
\label{tab:ae-results}
\renewcommand{\arraystretch}{1.25}
\setlength{\tabcolsep}{4pt}
\small
\begin{tabularx}{\linewidth}{@{}lXlXX@{}}
\toprule
\textbf{Источник} & \textbf{Архитектура} & \textbf{Датасет} &
\textbf{Постановка} & \textbf{Опубликованный результат} \\
\midrule
Mirsky и соавт.~\cite{mirsky2018kitsune} &
Kitsune (ансамбль AE с инкрементальным обучением) &
Mirai, OS Scan, SSL Renegotiation, SSDP Flood, ARP MitM, Active Wiretap, Fuzzing, Video Injection (собственный стенд) &
\eng{semi-supervised, online} &
По 9~атакам опубликованы значения AUC, среди которых:
Mirai $\approx 0{,}9988$, SSDP Flood $\approx 0{,}9844$, ARP MitM $\approx 0{,}9606$ \\
\addlinespace
Lopez-Martin и соавт.~\cite{lopez2017cvae} &
Conditional VAE &
NSL-KDD &
multiclass (4 класса) &
По данным авторов CVAE превосходит ряд классических классификаторов
(RF, MLP) на схеме оценки той же работы; конкретные значения см. в первоисточнике \\
\addlinespace
Andresini и соавт.~\cite{andresini2021autoencoder} &
AE-based deep metric learning &
NSL-KDD, AAGM, CICIDS2017 &
binary &
Авторы сообщают о превосходстве предложенного подхода над baseline-AE и
RF/SVM-классификаторами по F1 на трёх датасетах \\
\bottomrule
\end{tabularx}
\end{table}

Замечания по применимости AE/VAE.
\begin{itemize}
    \item Kitsune и его варианты эффективны в \emph{semi-supervised}-режиме
    (обучение только на нормальном трафике), что выгодно отличает их от
    supervised DL-моделей в условиях ограниченной разметки.
    \item Сравнение с supervised-моделями некорректно по принципу: AE
    решают задачу детекции аномалий, в то время как CNN-LSTM и аналоги в
    данной работе нацелены на бинарную/мультиклассовую классификацию с
    разметкой. В настоящей работе AE и VAE используются как контрольные
    semi-supervised baselines.
\end{itemize}

\subsection{Сводное сравнение архитектур по опубликованным результатам}

В \tabref{tab:summary-architectures} приведено сводное сопоставление
архитектур по совокупности качественных характеристик и опубликованных
данных.

\begin{table}[H]
\centering
\caption{Сводное сопоставление архитектур по опубликованным
NIDS-результатам и теоретическим характеристикам}
\label{tab:summary-architectures}
\renewcommand{\arraystretch}{1.25}
\setlength{\tabcolsep}{3pt}
\footnotesize
\begin{tabularx}{\linewidth}{@{}lXccXl@{}}
\toprule
\textbf{Архитектура} & \textbf{Локальные / временные структуры} &
\textbf{Стоимость обучения} & \textbf{Inference latency} &
\textbf{Опубликованные ориентиры качества (NIDS)} & \textbf{Источник} \\
\midrule
DNN/MLP & нет/нет        & низкая   & низкая   &
наибольшая Accuracy на NSL-KDD/UNSW-NB15 в семействе DNN-1\dots5 &
\cite{vinayakumar2019dl, sharafaldin2018toward, moustafa2016eval} \\
\addlinespace
1D-CNN  & да/огр.        & средняя  & низкая   &
устойчивое преимущество над DNN на тех же датасетах в одной и той же работе &
\cite{vinayakumar2019dl, ferrag2020dlcyber, liu2019survey} \\
\addlinespace
LSTM    & нет/да         & средняя  & средняя  &
\(\approx 96{,}93\%\) Acc на KDD Cup 99 в~\cite{kim2016lstm}; уступает
CNN-LSTM на UNSW-NB15 в~\cite{vinayakumar2019dl} &
\cite{kim2016lstm, yin2017dlrnn, vinayakumar2019dl} \\
\addlinespace
GRU     & нет/да         & средняя  & низкая   &
сопоставима с LSTM при меньшем числе параметров &
\cite{vinayakumar2019dl, ahmad2021nids} \\
\addlinespace
BiLSTM  & нет/двунаправ. & высокая  & средняя  &
прирост Accuracy/F1 по сравнению с LSTM на NSL-KDD и UNSW-NB15 &
\cite{ahmad2021nids, ferrag2020dlcyber} \\
\addlinespace
\textbf{CNN-LSTM} & \textbf{да/да} & \textbf{высокая} & \textbf{средняя} &
\textbf{превосходит DNN/CNN/LSTM в одной и той же работе на UNSW-NB15 и CICIDS2017} &
\textbf{\cite{vinayakumar2019dl, ferrag2020dlcyber, liu2019survey, ahmad2021nids}} \\
\addlinespace
TCN         & да/фиксиров. & средняя  & низкая   &
NIDS-публикаций существенно меньше; в обзорах рассматривается как
альтернатива LSTM при ограничениях по latency &
\cite{bai2018tcn, ferrag2020dlcyber} \\
\addlinespace
Transformer & через attention/да & высокая & средняя &
сопоставимое с CNN-LSTM качество в обзорах при повышенной дисперсии по seed &
\cite{vaswani2017attention, ferrag2020dlcyber, ahmad2021nids} \\
\addlinespace
AE          & косвенно/нет  & средняя & низкая &
ансамбль AE Kitsune: AUC $> 0{,}96$ для большинства атак &
\cite{mirsky2018kitsune, andresini2021autoencoder} \\
\addlinespace
VAE         & косвенно/нет  & средняя & низкая &
CVAE превосходит ряд классических baselines на NSL-KDD &
\cite{lopez2017cvae, kingma2014vae} \\
\bottomrule
\end{tabularx}
\end{table}

В \tabref{tab:summary-architectures} обоснованно выделена строка
CNN-LSTM: гибрид сочетает локальное (CNN) и временное (LSTM) моделирование,
а в нескольких независимых рецензируемых
источниках~\cite{vinayakumar2019dl, ferrag2020dlcyber, liu2019survey,
ahmad2021nids} устойчиво демонстрирует одно из лучших значений F1 и
Accuracy в внутренних сравнениях этих работ.

\subsection{Ограничения сравнения и угрозы валидности}

При интерпретации опубликованных результатов следует учитывать следующие
ограничения, систематически отмечаемые в
обзорной литературе~\cite{khraisat2019survey, sommer2010outside,
pendlebury2019tesseract, gharib2016evalframework}.

\begin{enumerate}
    \item \textbf{Различия датасетов.} KDD Cup 99 и NSL-KDD не отражают
    современную модель угроз; результаты на них завышают эффект
    архитектуры и не переносятся на эксплуатационные сценарии. Следовательно,
    высокие метрики LSTM-RNN на KDD Cup 99 (Kim и соавт.,
    \(\approx 96{,}93\%\) Accuracy~\cite{kim2016lstm}) нельзя
    интерпретировать как доказательство практической пригодности.
    \item \textbf{Отсутствие контроля дрейфа.} Большинство работ оценивают
    модели на статически фиксированных тестовых выборках, что приводит к
    \emph{spatial bias}~\cite{pendlebury2019tesseract}; реальная
    устойчивость моделей в потоковом сценарии остаётся неосвещённой.
    \item \textbf{Несбалансированный учёт ложных срабатываний.} Опубликованные
    значения часто относятся к Accuracy, не учитывающей стоимости FP при
    дисбалансе. Работы, приводящие FP-rate, FAR или FP-per-time, остаются
    в меньшинстве.
    \item \textbf{Несравнимость baselines.} Между работами различается состав
    baselines, что затрудняет ранжирование. Например,
    в~\cite{vinayakumar2019dl} и~\cite{ahmad2021nids} CNN-LSTM сравнивается
    с разными конфигурациями DNN/CNN/LSTM, и различие в относительных
    приростах не указывает на устойчивое превосходство.
\end{enumerate}

Эти ограничения определяют необходимость собственного эксперимента в равных
условиях, описание которого приведено в разделе~3 настоящего отчёта (Этап~2)
и реализуется на Этапе~4 ВКР.

\subsection{Обоснование выбора CNN--LSTM в качестве основной модели работы}
\label{subsec:why-cnnlstm}

С учётом теоретического анализа Этапа~1 и систематизированных опубликованных
результатов, представленных в \tabref{tab:summary-architectures}, в качестве
основной (\emph{proposed}) модели работы выбирается CNN--LSTM. Решение
обосновывается следующими аргументами.

\begin{enumerate}
    \item \textbf{Соответствие природе данных.} Атаки в flow-представлении
    проявляются одновременно как локальные паттерны (например, специфическое
    сочетание признаков SYN-флуда в одном flow) и как временные тренды
    (последовательность подобных flow). CNN--LSTM моделирует обе компоненты
    в одной архитектуре~\cite{vinayakumar2019dl, ferrag2020dlcyber}.
    \item \textbf{Опубликованные результаты.} В нескольких независимых
    рецензируемых работах~\cite{vinayakumar2019dl, ferrag2020dlcyber,
    liu2019survey, ahmad2021nids} CNN-LSTM демонстрирует наибольшие F1 и
    Accuracy в \emph{внутреннем} сравнении с DNN, CNN, LSTM и BiLSTM
    на UNSW-NB15 и CICIDS2017.
    \item \textbf{Воспроизводимость.} Архитектура хорошо документирована,
    реализуется стандартными примитивами PyTorch, не требует
    специализированных техник обучения (как Transformer на ограниченных
    выборках), что обеспечивает воспроизводимость экспериментов.
    \item \textbf{Эксплуатационная стоимость.} Inference latency
    CNN-LSTM, согласно обзорам~\cite{ferrag2020dlcyber, ahmad2021nids},
    укладывается в диапазон, приемлемый для NTA-pipeline (десятки
    миллисекунд при batch=1 на CPU); это совместимо с требованием NFR-3
    раздела~4 Этапа~1.
    \item \textbf{Альтернативы и контрольные модели.} В качестве альтернатив
    в собственном эксперименте сохраняются BiLSTM (двунаправленный контекст)
    и Transformer-encoder (attention); в качестве контрольных
    semi-supervised baselines --- AE и VAE; в качестве классических baseline
    --- Logistic Regression, Random Forest, XGBoost, Isolation Forest,
    One-Class SVM.
\end{enumerate}

\subsection{Промежуточные выводы по разделу}

\begin{enumerate}
    \item Систематизированы опубликованные результаты применения восьми
    типов архитектур в NIDS-задачах, эквивалентных постановке настоящей
    работы. Все результаты приведены с указанием источника, датасета и
    постановки.
    \item Зафиксированы три группы методологических ограничений сравнения
    (несовместимость датасетов, отсутствие контроля дрейфа, доминирование
    Accuracy при дисбалансе классов), что определяет необходимость
    собственного эксперимента в равных условиях.
    \item Сводное сопоставление по характеристикам архитектур и опубликованным
    NIDS-результатам подтверждает теоретический выбор CNN-LSTM как основной
    архитектуры; альтернативами в равных условиях останутся BiLSTM и
    Transformer-encoder, контрольными semi-supervised baselines --- AE и VAE,
    классическими baseline --- LR/RF/XGBoost/IF/OCSVM.
    \item Полученное обоснование является \emph{теоретическим} и подтверждается
    \emph{ссылочным} образом; \emph{количественное} подтверждение возможно
    только в собственном эксперименте, выполняемом на Этапе~4 ВКР по
    протоколу, формализованному в разделе~3 настоящего отчёта.
\end{enumerate}
